\chapter{Learning the Reactor Response and Enforcing Physical Admissibility}
\label{ch:projection}

\section{Accuracy and physical admissibility as distinct questions}

A surrogate learns the relationship between a reactor input and its steady-state effluent. Its statistical task is to reproduce component concentrations from examples. Its physical task is to return a nonnegative component vector with the invariant values appropriate to the supplied influent. These tasks are related because the mechanistic targets satisfy the reactor equations. They remain distinct because an ordinary regression loss does not explicitly restrict every prediction to the sample-specific feasible set \citep{Kashinath2021,Hansen2024}.

The complete input for sample $i$ is
\begin{equation}
x_i=(u_i^\top,c_{in,i}^\top)^\top\in\mathbb{R}^{22},\qquad
u_i=(\mathrm{HRT}_i,u_{\mathrm{air},i})^\top.
\label{eq:surrogate_input_vector}
\end{equation}
A fitted surrogate $f_m$ produces the raw component prediction
\begin{equation}
c_{raw,i}^{(m)}=f_m(x_i)\in\mathbb{R}^{20}.
\label{eq:surrogate_raw_state}
\end{equation}
The raw prediction is assessed in its returned physical units. A post prediction projection then produces a second prediction for the same sample. Pairing the two states permits the effect of physical enforcement to be assessed separately from the choice of learner, data partition, and preprocessing.

The comparison includes 13 statistical surrogates. The set spans boosting, randomized ensembles, kernel regression, neighborhood regression, latent-variable regression, sparse multiresponse regression, and neural networks. This breadth addresses the model-selection concern raised by \citet{BhosekarIerapetritou2018}. The surveys of \citet{Borisov2024} and \citet{ShwartzZiv2022} further explain why a neural architecture alone should not determine the comparator set for a tabular prediction problem. The purpose of the common protocol is to compare explicitly specified predictors for the same physical state.

\section{Constructing a mechanistic learning domain}

The single-reactor dataset is designed to contain 10,000 accepted steady states. Each candidate varies all 22 input coordinates. A randomized, scrambled, strength-one Latin hypercube places one sample in each stratum of every marginal input distribution. The construction uses seed 42 and no subsequent design optimization. Linear scaling maps the unit-hypercube samples to the bounds in Table~\ref{tab:reactor_sampling_domain}. Latin hypercube sampling distributes coverage across marginal ranges without requiring a full Cartesian grid \citep{McKay1979}. It does not guarantee uniform coverage of every nonlinear response feature or every high-dimensional interaction.

\begin{table}[htbp]
\centering\small
\caption{Design bounds for the standalone reactor learning domain. These intervals specify sampling limits and do not report realized sample extrema.}\label{tab:reactor_sampling_domain}
\begin{tabular}{lp{5.7cm}ll}
\toprule
Input & Meaning & Interval & Unit\\
\midrule
HRT & Hydraulic operating input & $[6,36]$ & h\\
$u_{\mathrm{air}}$ & Dimensionless aeration input & $[0.5,2.5]$ & 1\\
$S_O$ & Influent dissolved oxygen & $[0,0.5]$ & g O$_2$ m$^{-3}$\\
$S_F$ & Influent fermentable substrate & $[20,180]$ & g COD m$^{-3}$\\
$S_A$ & Influent acetate & $[5,80]$ & g COD m$^{-3}$\\
$S_{NH4}$ & Influent ammonium nitrogen & $[12,55]$ & g N m$^{-3}$\\
$S_{NO2}$ & Influent nitrite nitrogen & $[0,3]$ & g N m$^{-3}$\\
$S_{NO3}$ & Influent nitrate nitrogen & $[0,8]$ & g N m$^{-3}$\\
$S_{N2}$ & Influent dissolved dinitrogen & $[0,2]$ & g N m$^{-3}$\\
$S_{PO4}$ & Influent soluble phosphate & $[2,18]$ & g P m$^{-3}$\\
$S_I$ & Influent soluble inert organic matter & $[10,90]$ & g COD m$^{-3}$\\
$S_{ALK}$ & Influent alkalinity & $[80,260]$ & mol m$^{-3}$\\
$X_I$ & Influent particulate inert matter & $[20,120]$ & g COD m$^{-3}$\\
$X_S$ & Influent slowly biodegradable substrate & $[60,280]$ & g COD m$^{-3}$\\
$X_H$ & Influent heterotrophic biomass & $[15,100]$ & g COD m$^{-3}$\\
$X_{PAO}$ & Influent phosphate-accumulating biomass & $[5,60]$ & g COD m$^{-3}$\\
$X_{PP}$ & Influent stored polyphosphate & $[2,20]$ & g P m$^{-3}$\\
$X_{PHA}$ & Influent stored polyhydroxyalkanoates & $[1,30]$ & g COD m$^{-3}$\\
$X_{AOB}$ & Influent ammonia-oxidizing biomass & $[0.5,8]$ & g COD m$^{-3}$\\
$X_{NOB}$ & Influent nitrite-oxidizing biomass & $[0.5,8]$ & g COD m$^{-3}$\\
$X_{MeP}$ & Influent precipitated metal phosphate & $[0,12]$ & g TSS m$^{-3}$\\
$X_{MeOH}$ & Influent metal hydroxide & $[0,12]$ & g TSS m$^{-3}$\\
\bottomrule
\end{tabular}
\end{table}

The standalone domain is a declared simulation design. It is not a distribution estimated from plant measurements. In particular, its alkalinity interval is retained in the stated native coordinate. The operating and influent domains of a coupled plant are specified independently. A common component symbol does not establish that two evaluation designs use the same bounds.

For each candidate input, nonlinear least squares solves Equation~\eqref{eq:reactor_steady_balance} with numerical state bounds $[10^{-8},1500]$ in each component's native unit. The cost-change, step-size, and gradient convergence tolerances are $10^{-9}$. Each solve allows at most 1200 function evaluations. The generation procedure has an overall attempt cap of 25 times the target retained count. A state is retained only if the solver converges and the maximum absolute component-balance residual is no greater than $10^{-9}$ on the fixed native component-unit-per-day coordinates. The residual criterion is a numerical acceptance rule rather than a dimensionless physical error measure.

Two physically informed starting states are tried before dynamic relaxation. The first blends the sampled influent with a preceding accepted effluent when one is available. It lowers readily consumed substrate coordinates, retains biomass pools, and estimates dissolved oxygen from dilution and transfer. The second applies stronger substrate depletion and biomass enrichment. Starting vectors are restricted to $[10^{-6},1500]$. If both steady solves fail, a stiff backward differentiation formula integrates the dynamic balance for 20 days and supplies a new initial state. The state bounds and starting-state restrictions apply to concentrations. The 28 kinetic process rates are evaluated without an artificial ceiling or clipping.

Solver acceptance, invariant residual, and minimum concentration are recorded separately. A converged bounded least-squares calculation can terminate without meeting the required balance residual. A positive state can also fail conservation. Retaining these separate diagnostics prevents one successful numerical check from standing in for all others. The surrogate inputs contain the sampled operating and influent coordinates. Initialization choices and solver history are not predictive features.

\section{How the thirteen predictors represent the response}

\subsection{Boosting and randomized tree ensembles}

XGBoost builds an additive sequence of regression trees. Each new tree addresses the error left by the current ensemble while regularization discourages an unnecessarily complex fit. Its split structure represents changes in response across input regions, and repeated additions permit interactions and nonlinear response shapes. The computational design and regularized objective of \citet{ChenGuestrin2016} motivate its inclusion as a flexible tabular baseline.

LightGBM also uses gradient-boosted trees, but grows trees through leaf-wise expansion. This strategy can allocate more detail to regions where an additional split offers a larger loss reduction. The efficiency-oriented design of \citet{Ke2017} makes the method relevant when each candidate fit must be repeated over several validation partitions. Its learning rate, leaf count, and minimum leaf support jointly control how closely it follows the training response.

CatBoost uses symmetric boosted trees and a Bayesian bootstrap in the specified regression setting. The work of \citet{Prokhorenkova2018} addresses bias associated with categorical-feature handling and boosting. The reactor benchmark has continuous inputs, so its inclusion evaluates the corresponding regression learner without claiming a categorical-data advantage. Each level of a symmetric tree uses the same splitting condition across nodes at that level. This structural restriction gives a different approximation class from the leaf-wise trees of LightGBM.

Adaptive boosting (AdaBoost) regression combines shallow trees through adaptive emphasis on observations that are difficult to predict. The base trees have depth three. The estimator count, learning rate, and loss choice control how strongly the ensemble responds to poorly fitted examples. \citet{Drucker1997} develops the regression extension of boosting that supports this comparator. These four boosting methods fit one regressor for each effluent component. All 20 regressors within a method use one shared selected hyperparameter vector.

Random Forest averages randomized regression trees. Variation in sampled observations and candidate features reduces dependence among trees, so the average is less sensitive to a particular tree fit \citep{Breiman2001}. Extra Trees adds randomization of split thresholds \citep{Geurts2006}. Both methods fit the 20 outputs jointly in native multiresponse ensembles. Joint output handling lets the same tree partition support several component predictions. It does not impose a stoichiometric relationship among those outputs.

Tree averaging can preserve nonnegativity when all leaf responses are nonnegative. It cannot generally preserve the balance target for a new influent. For a hypothetical two-component system, two training responses with totals of 10 and 20 have an equally weighted average with total 15. If a new influent requires total 12, the average is positive but not conservative for that sample. This distinction is caused by the sample-specific target $b_i$, rather than by whether the predictor averages positive observations.

\subsection{Kernel, neighborhood, and latent-variable regression}

Support vector regression fits a function with an error-insensitive region around each training response. Predictions within that region incur no corresponding error penalty. A linear kernel gives a globally linear input mapping. A radial kernel allows a nonlinear mapping through similarity to training inputs. The penalty parameter balances fit and regularity, while the insensitive width sets the tolerated response discrepancy. \citet{SmolaScholkopf2004} explains these mechanisms and the importance of their parameterization. The benchmark fits one regressor per component with one shared setting vector.

The $k$-nearest-neighbor predictor finds training observations close to a new input and combines their 20-component responses. The neighborhood size controls the degree of smoothing. Uniform weights give every selected neighbor equal influence, whereas distance weights give closer observations greater influence. The local behavior analyzed by \citet{Mack1981} explains why input scaling and distance choice matter. An unscaled alkalinity coordinate spanning hundreds of numerical units could dominate a much smaller aeration coordinate even though both influence the reactor. Standardization avoids using native numerical range as an implicit measure of scientific importance.

Partial least squares regression represents the inputs and joint response through latent directions chosen to capture their covariance. It can summarize strongly related input variables with a relatively small number of components. The chemometric formulation described by \citet{Wold2001} provides an inspectable linear baseline. Its wastewater use is illustrated by \citet{Woo2009}, who employs a kernel extension for process-variable estimation, and by \citet{Yang2021}, who combines partial least squares with another predictor for adaptive effluent-quality prediction. The benchmark retains ordinary partial least squares with one to six latent components. The nonlinear extensions in those applications do not make the present baseline nonlinear.

\subsection{Shared sparse coefficient models}

Multi-task Lasso fits one coefficient matrix for all component outputs. If row $B_{j,\cdot}$ contains the coefficients of input $j$ across those outputs, the row-group penalty is
\begin{equation}
\|B\|_{2,1}=\sum_j\|B_{j,\cdot}\|_2.
\label{eq:multiresponse_row_penalty}
\end{equation}
This penalty can set an entire input row to zero. It therefore promotes a shared selection of inputs rather than 20 unrelated selections. The feature-sharing principle of \citet{Argyriou2008}, the multi-task Lasso algorithm of \citet{Liu2009}, and the grouped-variable framework of \citet{Yuan2006} provide the relevant foundations.

Multi-task Elastic Net combines the row-group penalty with a squared Frobenius penalty. The row-group part encourages selection and the squared part shrinks the retained coefficients. This extends the elastic-net principle of \citet{Zou2005} to a joint response structure, with multiresponse computation described by \citet{Simon2013}. Both methods use standardized inputs and targets. Their global coefficient matrices can be inspected to identify fitted associations. A nonzero row shows that an input contributes to the fitted prediction, but it does not establish a causal effect or a mechanistic reaction rate.

These two predictors and partial least squares are grouped as structurally inspectable surrogates. The grouping identifies their globally available coefficient structures. It is not a measured claim that users understand them more easily or that their coefficients are scientifically correct. It also does not establish physical admissibility. A global linear relation can still predict a negative concentration or violate the influent-specific balance.

\subsection{Neural predictors}

The multilayer perceptron uses three fully connected hidden layers with 128, 64, and 32 units and one joint 20-output layer. Each hidden layer combines its inputs and applies a nonlinear activation. Backpropagation updates the weights according to the prediction loss \citep{Rumelhart1986}. The approximation result of \citet{Hornik1989} motivates the expressive capacity of feedforward networks. It does not guarantee that this finite architecture learns the reactor response from a given sample size. The topology considerations discussed by \citet{Curteanu2011} likewise motivate explicit architectural and training limits rather than an unspecified neural comparator.

TabNet applies sequential attentive feature selection. Different decision steps can emphasize different subsets of input features before their contributions enter a joint output layer. \citet{ArikPfister2021} develops this approach for tabular learning. The attention masks provide information about the model's use of inputs, but they do not impose conservation or make the component outputs causal explanations. TabNet remains a neural comparator even though its internal masks can be inspected.

Both neural predictors start from random initialization and use only the simulated training partition available to that fit. No external learned representation is used. The multilayer perceptron uses all available training rows and stops on training-loss convergence or at 250 epochs. Validation-based early stopping is disabled. Its training-loss patience is 15 epochs. TabNet's epoch count is selected in the inner search and ranges from 20 to 50. Its patience is zero, so a refit runs the selected epoch count without reserving an additional response-based stopping partition.

\section{Preprocessing and validation without response leakage}

The evaluation uses nested subsets at
\begin{equation}
N\in\{500,1450,2400,3350,4300,5250,6200,7150,8100,9050,10000\}.
\label{eq:reactor_nested_sample_sizes}
\end{equation}
A seed-42 permutation determines the nested row order and persistent master fold identifiers before subsets are formed. Every smaller subset is contained in each larger subset. A row retains its outer-fold role at every scale. Five outer folds train on $0.8N$ observations and score the remaining $0.2N$. The same rows and fold identifiers are used by every surrogate and both prediction types.

Hyperparameters are selected using four inner folds within the largest outer-training partition. The scored outer fold is absent from the search. For each outer fold, its selected setting is retained over the smaller nested subsets. Learning curves therefore examine the response to sample size under a fixed fold-specific setting. They do not represent a new optimum obtained through a separate hyperparameter search at every scarce-data scale. The bias--variance discussion of \citet{Belkin2019} cautions against assuming that error must vary monotonically with model capacity or available data. The learning-curve review of \citet{VieringLoog2023} supports treating the complete curve as an object of assessment.

Every transformation is fitted within the applicable training partition. During inner selection, that partition is the inner-training split. During outer scoring, it is the outer-training subset at the current size. During extrapolation assessment, it is the full in-domain training pool. The mean and population standard deviation of a training coordinate define its z-score transformation. Validation inputs are transformed with those training quantities, and target-scaled predictions are restored to physical units before projection or physical diagnostics.

\begin{table}[htbp]
\centering
\caption{Output handling and preprocessing for the statistical comparators. Scaling uses only the training partition applicable to each fit.}
\label{tab:surrogate_preprocessing}
\small
\begin{tabular}{p{5.1cm}p{3.2cm}p{3.3cm}}
\toprule
Surrogates & Output handling & Scaling\\
\midrule
XGBoost, LightGBM, CatBoost & Separate component regressors & None\\
AdaBoost & Separate component regressors & Targets only\\
Random Forest & Joint tree response & None\\
Extra Trees & Joint tree response & Targets only\\
Support vector regression & Separate component regressors & Inputs and targets\\
$k$-nearest neighbors & Joint neighborhood response & Inputs and targets\\
Partial least squares, Multi-task Lasso, Multi-task Elastic Net & Joint coefficient structure & Inputs and targets\\
Multilayer perceptron, TabNet & Joint neural output & Inputs and targets\\
\bottomrule
\end{tabular}
\end{table}

Tree partitions depend on coordinate ordering, so a positive affine input rescaling preserves candidate threshold partitions. Distance-based, covariance-based, and regularized models require more deliberate scaling because relative numerical magnitudes affect their objectives. Separate-output tree regressors do not combine different component units within one joint fitting loss. The target-only scaling of Extra Trees and AdaBoost is retained as a specified benchmark choice. Independent of those training choices, every scored component error is standardized with the designated training reference. Training target scaling and evaluation normalization are distinct operations.

No missing-value imputation, response-based outlier removal, or target filtering is applied. No validation or extrapolation response is used to estimate a scaling parameter. A hypothetical decision to standardize a validation target with the standard deviation of its own fold would change the scoring reference using data outside training. Retaining a training reference instead asks how large an error is relative to the variability the learner was allowed to observe.

\section{Hyperparameter selection and finite comparison scope}

Each surrogate uses a seeded tree-structured Parzen search with 100 trials per study and no pruning or timeout. The search mechanism is supported by the optimization framework of \citet{Akiba2019}. The objective is the mean squared standardized component error over the inner validation folds. Five studies select the settings for the five largest-scale outer-training partitions. A sixth, independent four-fold study selects the setting for a final full-domain fit. The design specifies 600 trials per surrogate and 7800 across all 13 surrogates. These are allocated search budgets rather than performance findings.

\begin{table}[htbp]
\centering\small
\caption{Search domains and fixed controls for tree ensemble surrogates. Integer limits are inclusive. ``Log'' denotes a log-uniform search.}\label{tab:surrogate_search_domains}
\begin{tabular}{p{2.65cm}p{9.2cm}}
\toprule
Surrogate & Search domain and fixed controls\\
\midrule
XGBoost & 80--320 trees. Learning rate 0.01--0.2 (log). Depth 3--8. Minimum child weight 0.5--8 (log). Row and feature fractions 0.6--1. Absolute-value penalty $10^{-6}$--1 (log). Squared penalty $10^{-3}$--10 (log). Histogram or approximate trees. Squared-error objective.\\
LightGBM & 80--320 trees. Learning rate 0.01--0.2 (log). 15--63 leaves. Depth unbounded or in $\{4,6,8,10\}$. Minimum child rows 5--40. Row and feature fractions 0.6--1 with subsampling frequency 1. Absolute-value penalty $10^{-6}$--1 (log). Squared penalty $10^{-6}$--10 (log). Regression objective.\\
CatBoost & 80--320 iterations. Learning rate 0.01--0.2 (log). Depth 4--8. Squared penalty 0.01--20 (log). Bagging temperature 0--5. Random strength $10^{-3}$--5 (log). Root mean squared error loss and Bayesian bootstrap.\\
AdaBoost & 50--250 depth-three trees. Learning rate 0.01--1 (log). Linear, squared, or exponential regression loss.\\
Random Forest & 100--400 trees. Depth unbounded or in $\{6,10,14,18\}$. Minimum split size 2--10. Minimum leaf size 1--4. Feature fraction 0.5--1. Bootstrap enabled or disabled.\\
Extra Trees & 100--400 trees. Depth unbounded or in $\{6,10,14,18\}$. Minimum split size 2--10. Minimum leaf size 1--4. Feature fraction 0.5--1. Bootstrap enabled or disabled.\\
\bottomrule
\end{tabular}
\end{table}

\begin{table}[htbp]
\centering\small
\caption{Search domains and fixed controls for kernel, neighbor, multiresponse, and neural surrogates. Integer limits are inclusive. ``Log'' denotes a log-uniform search.}\label{tab:surrogate_search_domains_other}
\begin{tabular}{p{2.65cm}p{9.2cm}}
\toprule
Surrogate & Search domain and fixed controls\\
\midrule
Support vector regression & Penalty 0.1--10 (log). Error-insensitive width 0.01--0.5 (log). Linear or radial kernel. Radial coefficient set from standardized input variance or inverse input dimension. Tolerance $10^{-3}$--$10^{-2}$ (log). Shrinking enabled. At most 20,000 iterations. Kernel cache of 256 megabytes per output regressor.\\
$k$-nearest neighbors & 1--15 neighbors. Uniform or distance weights. Automatic, ball-tree, $k$-dimensional tree, or brute-force search. Leaf size 15--60. Minkowski, Euclidean, or Manhattan distance. Minkowski order in $\{1,2,3\}$.\\
Partial least squares & 1--6 latent components. 200--1000 maximum iterations. Tolerance $10^{-8}$--$10^{-4}$ (log). Internal scaling disabled because external training-partition scaling is applied.\\
Multi-task Elastic Net & Penalty strength $10^{-6}$--10 (log). Row-group mixing ratio 0.05--0.95. Cyclic coordinate descent. At most 20,000 iterations and tolerance $10^{-6}$.\\
Multi-task Lasso & Penalty strength $10^{-6}$--10 (log). Cyclic coordinate descent. At most 20,000 iterations and tolerance $10^{-6}$.\\
Multilayer perceptron & Hidden widths $(128,64,32)$. Rectified linear, hyperbolic tangent, or logistic activation. Adaptive moment optimization. Squared-weight penalty $10^{-6}$--$10^{-2}$ (log). Batch size in $\{64,128,256\}$. Learning rate $5\times10^{-4}$--$10^{-2}$ (log). Training-loss tolerance $10^{-4}$--$10^{-3}$ (log). Shuffling enabled. Training-loss patience 15 epochs. At most 250 epochs.\\
TabNet & Equal decision and attention widths in $\{8,16,24\}$. Three or four attention steps. Relaxation coefficient 1--2. Sparsity coefficient $10^{-6}$--$10^{-2}$ (log). Adaptive moment learning rate 0.003--0.03 (log). Sparsemax or entmax-1.5 masks. 20--50 selected epochs. Batch size in $\{512,1024\}$. Virtual batch size 128. Patience zero.\\
\bottomrule
\end{tabular}
\end{table}

The finite ranges define the actual comparison. A linear or radial support vector model does not represent every possible kernel configuration. A neural predictor limited to 250 epochs does not represent every training budget or architecture. The shared 100-trial allocation does not guarantee equal difficulty of tuning across model families. The study therefore assesses these specified approximation classes under a common protocol rather than claiming universal superiority of one family.

\section{The Kircher--Votsmeier conservative projection}

\subsection{A positive provisional target and relative-error geometry}

\citet{KircherVotsmeier2025} develop a correction that combines conservation with positivity for mechanistic-kinetics surrogates. Its transfer to the reactor prediction problem uses the invariant operator $A$, the basis $Z$, and the feasible influent anchor. The projection depends on the completed component prediction rather than on predictor internals. It can therefore be applied to every comparator without changing its training objective or architecture. It still requires the full component vector, the corresponding nonnegative influent, and the physical operators.

The first step forms a strictly positive provisional target,
\begin{equation}
c_{+,i}^{(m)}=\max(c_{raw,i}^{(m)},\boldsymbol{\varepsilon}_c),\qquad
W_i^{(m)}=\operatorname{Diag}\bigl(\mathbf{1}\oslash c_{+,i}^{(m)}\bigr).
\label{eq:projection_positive_target}
\end{equation}
The maximum and division are componentwise. Each entry of $\boldsymbol{\varepsilon}_c$ is $10^{-8}$ in its component's native unit. The provisional target supplies both a reference state and finite positive weights. It is an intermediate construction, so its clipped coordinates are not the final physical prediction.

For a component change $h_j$, the weighted discrepancy is $h_j/c_{+,ij}^{(m)}$. The geometry therefore penalizes relative movement. A hypothetical movement of 1 g m$^{-3}$ corresponds to a relative discrepancy of 0.01 at a component value of 100 g m$^{-3}$ and 10 at a component value of 0.1 g m$^{-3}$. Small coordinates receive strong protection. Large coordinates can absorb larger absolute changes at the same relative cost. This explains why relative-error correction and variance-standardized accuracy need not favor the same redistribution.

\subsection{Selecting an invariant-consistent candidate}

All invariant-consistent states for sample $i$ have the form
\begin{equation}
c=c_{in,i}+Zq,\qquad q\in\mathbb{R}^{15}.
\label{eq:projection_affine_parameterization}
\end{equation}
The projection selects the conservative change closest to the provisional target change in the relative-error geometry,
\begin{equation}
q_i^{*,(m)}=\underset{q\in\mathbb{R}^{15}}{\operatorname{argmin}}
\left\|W_i^{(m)}\bigl[Zq-(c_{+,i}^{(m)}-c_{in,i})\bigr]\right\|_2^2.
\label{eq:projection_weighted_least_squares}
\end{equation}
The squared objective is a quadratic function of $q$. Differentiation gives
\begin{equation}
Z^\top(W_i^{(m)})^2Zq_i^{*,(m)}
=Z^\top(W_i^{(m)})^2(c_{+,i}^{(m)}-c_{in,i}).
\label{eq:projection_normal_conditions}
\end{equation}
For any nonzero vector $v$, the associated quadratic form equals $\|W_i^{(m)}Zv\|_2^2>0$. The matrix is therefore positive definite and the minimizer is unique.

The conservative candidate is
\begin{equation}
c_{cand,i}^{(m)}=c_{in,i}+Zq_i^{*,(m)}.
\label{eq:projection_conservative_candidate}
\end{equation}
Since $AZ=0$, it satisfies $Ac_{cand,i}^{(m)}=Ac_{in,i}=b_i$ in exact arithmetic. If the provisional target is already invariant-consistent, its change lies in $\operatorname{null}(A)$. The least-squares residual can then be zero and the candidate reproduces that target to numerical tolerance.

Equation~\eqref{eq:projection_normal_conditions} provides the derivation rather than the preferred numerical solution. The weighted design $W_i^{(m)}Z$ is solved through an orthogonal--triangular factorization. Singular-value decomposition with relative cutoff $10^{-12}$ is the fallback. No explicit inverse of the normal-equation matrix is formed. Near-zero provisional targets can produce very unequal weights, so solving the original weighted system avoids the extra conditioning penalty associated with forming a normal-equation inverse.

\subsection{The largest feasible interpolation toward the candidate}

Conservation of the candidate does not establish its nonnegativity. Define
\begin{equation}
\delta_i^{(m)}=c_{cand,i}^{(m)}-c_{in,i}=Zq_i^{*,(m)}.
\label{eq:projection_candidate_direction}
\end{equation}
The segment $c_{in,i}+\alpha\delta_i^{(m)}$, with $0\leq\alpha\leq1$, is conservative because $A\delta_i^{(m)}=0$. For every decreasing coordinate, nonnegativity requires
\begin{equation}
\alpha\leq\frac{c_{in,ij}}{-\delta_{ij}^{(m)}}
\qquad\text{when }\delta_{ij}^{(m)}<0.
\label{eq:projection_coordinate_bound}
\end{equation}
The minimum of these bounds is the first point at which the segment reaches the nonnegative-orthant boundary.

With numerical safety margin $\eta=10^{-12}$, the accepted interpolation factor and final state are
\begin{equation}
\begin{split}
\alpha_i^{(m)}&=
\begin{cases}
1, & c_{cand,i}^{(m)}\geq0,\\
(1-\eta)\displaystyle\min_{j\,\mid\,\delta_{ij}^{(m)}<0}
\frac{c_{in,ij}}{-\delta_{ij}^{(m)}}, & \text{otherwise},
\end{cases}\\
\tilde c_i^{(m)}&=c_{in,i}+\alpha_i^{(m)}\delta_i^{(m)}.
\end{split}
\label{eq:projection_positivity_backtracking}
\end{equation}
If the candidate is negative somewhere, at least one bound lies below one. Nonnegativity of the influent makes every bound nonnegative. The factor is restricted to $[0,1]$ against roundoff. Decreasing coordinates satisfy $\tilde c_{ij}^{(m)}\geq\eta c_{in,ij}\geq0$, and increasing coordinates remain nonnegative. Conservation follows from
\begin{equation}
A\tilde c_i^{(m)}=Ac_{in,i}+\alpha_i^{(m)}AZq_i^{*,(m)}=b_i.
\label{eq:projection_final_conservation}
\end{equation}
The unrestricted boundary factor is the largest feasible interpolation factor. The safety margin makes the returned factor slightly smaller when correction is necessary.

A hypothetical three-component system with conserved sum 10 illustrates this safeguard. Let the feasible influent be $(8,1,1)^\top$ and a conservative candidate be $(4,-0.5,6.5)^\top$. The direction is $(-4,-1.5,5.5)^\top$. Its decreasing-coordinate bounds are 2 and $2/3$, so the largest feasible factor is $2/3$. Ignoring the numerical safety margin only for this illustration gives $(16/3,0,14/3)^\top$. The sum remains 10 and every component is nonnegative. The first coordinate changes even though the second coordinate determines the boundary. One scalar shortens the complete conservative direction.

If an influent coordinate is zero and the candidate direction decreases that coordinate, its bound is zero. The final state then returns to the influent anchor. This is a valid feasible fallback, but it can remove useful information from the raw prediction. Likewise, a component that is absent from the invariant rows can still move during interpolation because the complete direction is shortened. Recording the factor and activation frequency is therefore necessary for interpreting the correction.

\subsection{Acceptance checks and limits of the guarantee}

A returned projection is accepted only when
\begin{equation}
\|A\tilde c_i^{(m)}-b_i\|_2\leq\tau_{mc},\qquad
\tilde c_{ij}^{(m)}\geq-\tau_{nn,j}\quad\text{for every }j,
\label{eq:projection_numerical_acceptance}
\end{equation}
with $\tau_{mc}=10^{-10}$ in the fixed invariant-coordinate convention and $\tau_{nn,j}=10^{-10}$ in the native unit of component $j$. These checks distinguish the exact mathematical construction from its finite-precision realization. A finite raw vector is required. The method does not define an admissible correction of a nonfinite prediction through this algebra.

The final state is feasible for the declared invariants and sign restrictions. It need not minimize the weighted distance over the full intersection $\mathcal{F}_i$. The least-squares stage minimizes distance over an affine space, and the interpolation stage searches one line segment. A general constrained nearest-point problem searches a larger set. The construction also does not minimize the variance-standardized prediction error, whose denominator differs from $c_+$. A reduction in conservation error therefore does not imply a reduction in statistical error.

No kinetic, equilibrium, stability, or field-performance guarantee follows from the projection. A conservative nonnegative state can have a large residual in Equation~\eqref{eq:reactor_steady_balance}. The distinction between physical reconciliation and generalization under a changed input distribution is consistent with \citet{Shen2023} and \citet{Valente2025}. The projection also assumes that the influent anchor is feasible for the correct boundary. An unmodeled external source invalidates that assumption even if the numerical correction succeeds.

Raw and projected composites are calculated with the same map,
\begin{equation}
y_{raw,i}^{(m)}=I_{comp}c_{raw,i}^{(m)},\qquad
\tilde y_i^{(m)}=I_{comp}\tilde c_i^{(m)}.
\label{eq:projection_paired_composites}
\end{equation}
Feasibility is checked before this aggregation. Composite values cannot reveal every negative coordinate or every inconsistent component redistribution.

\section{Evaluation of error, feasibility, and correction size}

\subsection{Component errors on a training-defined scale}

Let $\mathcal{T}$ denote the applicable training reference and let $\sigma_j^{(\mathcal{T})}$ be the population standard deviation of target component $j$ in that reference. For prediction type $t$, which is either raw or projected, define
\begin{equation}
\xi_{ij}^{(m,t)}=
\frac{c_{ij}^{(m,t)}-c_{out,ij}}{\sigma_j^{(\mathcal{T})}}.
\label{eq:projection_standardized_error}
\end{equation}
The normalized mean squared error (nMSE), normalized root mean squared error (nRMSE), and normalized mean absolute error (nMAE) are
\begin{align}
\mathrm{nMSE}_{m,t}&=\frac{1}{20n}\sum_{i=1}^{n}\sum_{j=1}^{20}(\xi_{ij}^{(m,t)})^2,\label{eq:projection_nmse}\\
\mathrm{nRMSE}_{m,t}&=\sqrt{\mathrm{nMSE}_{m,t}},\label{eq:projection_nrmse}\\
\mathrm{nMAE}_{m,t}&=\frac{1}{20n}\sum_{i=1}^{n}\sum_{j=1}^{20}|\xi_{ij}^{(m,t)}|.
\label{eq:projection_nmae}
\end{align}
The component state is the primary scoring space because it is the prediction target and the space in which constraints are imposed. The root mean squared score emphasizes larger standardized errors. The absolute score gives a complementary average deviation. A component with zero training standard deviation makes these normalized measures undefined and invalidates the run.

A hypothetical error of 1 g m$^{-3}$ has standardized magnitude 0.1 for a component with training standard deviation 10 g m$^{-3}$. The same physical error has magnitude 10 for a component with training standard deviation 0.1 g m$^{-3}$. The score gives equal numerical weight to standardized component variation. It does not claim that every component has equal regulatory importance or equal measurement uncertainty.

For a scored partition with true mean $\bar c_{out,j}$, the component coefficient of determination and its equal-weight average are
\begin{equation}
R_{j,m,t}^2=1-
\frac{\sum_i(c_{ij}^{(m,t)}-c_{out,ij})^2}
{\sum_i(c_{out,ij}-\bar c_{out,j})^2},\qquad
\bar R_{m,t}^2=\frac{1}{20}\sum_{j=1}^{20}R_{j,m,t}^2.
\label{eq:projection_r_squared}
\end{equation}
An undefined component denominator makes the corresponding macro-average undefined. The coefficient can be negative when the prediction is worse than the scored-partition mean. It complements the standardized error but uses a different response-variability reference.

Physical-unit errors are retained for individual components. Secondary composite error is calculated separately for COD, TN, TP, and TSS. The normalized root mean squared error of composite $q$ is its physical-unit root mean squared error divided by $\sigma_q^{(\mathcal{T})}$, the training standard deviation of that composite. This is not the primary component-average nRMSE. Composite standard deviations are calculated after applying the common reporting map to the training targets.

\subsection{Violation frequencies and their magnitudes}

The conservation residual and standardized total negative magnitude of a state $c$ are
\begin{equation}
v_{mc,i}(c)=\|Ac-b_i\|_2,\qquad
v_{nn,i}^{\mathrm{std}}(c)=
\left\|\min(c\oslash\boldsymbol{\sigma}^{(\mathcal{T})},0)\right\|_1.
\label{eq:projection_violation_magnitudes}
\end{equation}
The minimum is componentwise. The conservation residual uses the same full-precision orthonormal basis for every sample. It is a numerical norm of the adopted invariant coordinates rather than a direct mass concentration in one constituent unit. The negative magnitude uses training standard deviations to compare different component bases without summing incompatible physical units.

The two row-level violation frequencies are
\begin{equation}
\pi_{mc}=
\frac{1}{n}\sum_i\mathbb{1}[v_{mc,i}(c_i)>\tau_{mc}],\qquad
\pi_{nn}=
\frac{1}{n}\sum_i\mathbb{1}[\min_j(c_{ij}+\tau_{nn,j})<0].
\label{eq:projection_violation_rates}
\end{equation}
A state is counted as nonnegative only when every component passes its tolerance. Component-specific violation frequencies identify which coordinates cross the boundary. The joint feasible fraction is
\begin{equation}
\pi_{\mathrm{feasible}}=
\frac{1}{n}\sum_i\mathbb{1}
[v_{mc,i}(c_i)\leq\tau_{mc}\ \text{and}\ 
\min_j(c_{ij}+\tau_{nn,j})\geq0].
\label{eq:projection_feasible_fraction}
\end{equation}
Conservation and nonnegativity frequencies are not added to calculate this fraction because the same state can violate both requirements.

Mean, median, and maximum violation magnitudes are retained along with the frequencies. A hypothetical set of many concentrations just below zero can have a high violation frequency but small total negative magnitude. Another set with one strongly negative state can have low frequency and a large maximum magnitude. The two summaries answer different questions. Physical-unit negative magnitudes are also retained separately for each component.

\subsection{Paired accuracy effects and displacement}

For each outer fold, the effect on the primary error score is
\begin{equation}
\Delta\mathrm{nRMSE}_{m,N,g}=
\mathrm{nRMSE}_{m,\mathrm{projected},N,g}
-\mathrm{nRMSE}_{m,\mathrm{raw},N,g}.
\label{eq:projection_paired_error_difference}
\end{equation}
A negative difference indicates lower error after projection, and a positive difference indicates higher error. Differences are calculated within paired folds before summary. Component and composite differences are retained because a single aggregate score can conceal opposite changes among outputs.

The sample-level standardized displacement is
\begin{equation}
d_{c,i}^{\mathrm{std},(m)}=
\left\|(\tilde c_i^{(m)}-c_{raw,i}^{(m)})
\oslash\boldsymbol{\sigma}^{(\mathcal{T})}\right\|_2.
\label{eq:projection_standardized_displacement}
\end{equation}
This norm measures how far the correction moves the state relative to training variability. It measures movement rather than accuracy. A large displacement can repair a large physical inconsistency while moving away from the true kinetic response. Mean, median, and maximum displacement, componentwise physical-unit movement, the fraction requiring backtracking, and the distribution of $\alpha$ are therefore assessed together.

Learning curves retain the full error trajectory over Equation~\eqref{eq:reactor_nested_sample_sizes}. As a descriptive summary, the normalized area under an error curve is calculated by trapezoidal integration,
\begin{equation}
\mathcal{A}_{m,t}=
\frac{1}{N_{11}-N_1}
\sum_{\ell=1}^{10}\frac{E_{m,t}(N_\ell)+E_{m,t}(N_{\ell+1})}{2}
(N_{\ell+1}-N_\ell),
\label{eq:projection_learning_curve_area}
\end{equation}
where $E_{m,t}(N)$ is the mean held-out error at size $N$. This quantity has the same scale as the selected error measure. It summarizes behavior across the declared size interval and does not replace the individual size scores. Training and validation error at the largest size provide a separate descriptive gap. Neither summary establishes an optimal future sample size without an additional learning-curve model \citep{VieringLoog2023}.

Each in-domain metric is summarized as the arithmetic mean and sample standard deviation of five outer-fold values. The standard deviation uses denominator four and describes partition-to-partition variation. No hypothesis test or inferential model ranking follows from these summaries. They also do not supply calibrated predictive uncertainty for a new state. The ensemble uncertainty approach of \citet{Lakshminarayanan2017} addresses a different task, and no such uncertainty calibration is assumed here.

\section{Testing inputs beyond the learning domain}

Extrapolation assessment uses separately simulated inputs outside the training hyperrectangle. The final setting for each surrogate is selected with four folds on all 10,000 in-domain states. Its preprocessing and predictor are then refitted on that complete pool. No extrapolation response enters selection, fitting, scaling, or calibration. This separation addresses the domain-shift concern discussed in scientific machine learning by \citet{Karniadakis2021} and \citet{Shen2023}.

Four inputs can exceed their upper learning bounds. For input $\ell$ with bounds $[L_\ell,U_\ell]$, define its normalized exceedance as
\begin{equation}
e_{i\ell}=\frac{x_{i\ell}-U_\ell}{U_\ell-L_\ell}.
\label{eq:projection_domain_exceedance}
\end{equation}
The mild design uses $e_{i\ell}\in[0.10,0.25]$ and the severe design uses $e_{i\ell}\in[0.35,0.60]$. The generator targets a mixture of 70\% one-input excursions and 30\% two-input excursions. All remaining inputs remain within the original bounds. Every candidate therefore lies outside the training hyperrectangle and consequently outside the training-point convex hull.

\begin{table}[htbp]
\centering
\caption{Fixed exterior intervals for the two extrapolation designs. The intervals specify input conditions rather than prediction outcomes.}
\label{tab:reactor_extrapolation_domain}
\small
\begin{tabular}{lllll}
\toprule
Input & Learning interval & Mild interval & Severe interval & Unit\\
\midrule
HRT & $[6,36]$ & $[39,43.5]$ & $[46.5,54]$ & h\\
$u_{\mathrm{air}}$ & $[0.5,2.5]$ & $[2.7,3.0]$ & $[3.2,3.7]$ & 1\\
$S_{NH4}$ & $[12,55]$ & $[59.3,65.8]$ & $[70.1,80.8]$ & g N m$^{-3}$\\
$X_S$ & $[60,280]$ & $[302,335]$ & $[357,412]$ & g COD m$^{-3}$\\
\bottomrule
\end{tabular}
\end{table}

Each severity design targets 600 accepted steady states and retains the same solver, initialization sequence, and balance-residual criterion. Candidate, accepted, failed, valid-but-unselected, and retained counts are distinguished. The realized one-input and two-input proportions are also distinguished from their design targets. The intervals alone do not establish which exterior conditions are hardest to predict because reactor kinetics can respond differently to the four inputs.

Extrapolation errors are standardized with population component standard deviations from all 10,000 in-domain targets. Composite normalization uses the corresponding full-domain composite standard deviations. The nMSE, nRMSE, nMAE, conservation and nonnegativity diagnostics, displacement, and backtracking diagnostics are calculated on identical raw and projected cases. The coefficient of determination is secondary because the deliberately shifted response distribution changes its denominator.

Each exterior metric is a single descriptive summary over 600 cases from one full-domain refit. It is not an average over the five outer folds and has no fold standard deviation. Satisfaction of the projected constraints is interpreted separately from the predictive accuracy of extrapolation. A state can remain nonnegative and conservative while its estimated nitrogen distribution or biomass abundance is inaccurate.

\section{Computational cost and the meaning of a fast prediction}

Computation is partitioned into hyperparameter search, preprocessing and fitting, raw inference, projection alone, and complete projected inference. Search time is reported once for each distinct study because a largest-scale outer setting is reused at all eleven sizes. Summing that same search duration over the nested size rows would count the selection cost repeatedly. Fitting time excludes the search and reflects the preparation required for the specified training partition.

Inference uses a deterministic batch of 512 rows constructed by cycling through the relevant validation inputs. The batch size stays fixed even when fewer than 512 unique validation observations are available. Five warm-up calls precede 30 timed repetitions for each path. The median batch latency is divided by 512 and expressed as milliseconds per sample. For in-domain assessment, the five fold medians are summarized by their arithmetic mean and sample standard deviation. Exterior timing uses the same batching and repetitions for one final fit and has no fold uncertainty.

Projection timing includes positive-target construction, relative weighting, the weighted solve, conservation and nonnegativity checks, and any positivity backtracking. Complete projected inference includes both raw prediction and correction. All physical operators, returned component states, preprocessing, and metrics use double precision. TabNet's internal tensors use single precision before its predictions return to the common physical-coordinate calculations. Pseudorandom generation and stochastic fitting use streams derived from seed 42. Deterministic coordinate-descent settings are retained for the multi-task models.

These timing measures describe the declared batch protocol and computational environment. Dividing a batch duration by its row count does not establish the latency of a single isolated query. Likewise, low projection latency would not establish field readiness. A useful predictor must be assessed through its component accuracy, physical diagnostics, domain of use, preparation cost, and the behavior of downstream calculations. The post prediction correction provides a defined admissibility layer within that assessment.
